\section{Lessons for Practitioners}

\paragraph{Benchmark accuracy is the wrong selection criterion.} The most accurate classifier on the
benchmark was the least accurate a year later and on a new repository. Evaluate a labeler on issues
from after its training and calibration data, and on a held-out repository, before choosing it.

\paragraph{A calibrated threshold is not an error guarantee.} Tuning the threshold to the target on a
few hundred calibration issues over-promised in every condition. A risk-controlling procedure needs
its threshold candidates fixed before certification, and with 330 certification issues it certified
a 10\% target only for the classical classifier and at low coverage. Choose the target from a
frontier like Table~\ref{tab:frontier}, or budget for far more labeled calibration issues. A
certificate holds only for data exchangeable with the calibration set: recalibrate, on that
repository's own labeled issues, when a repository is onboarded.

\paragraph{Calibration across repositories fixes sets, not thresholds.} Leave-one-repository-out
calibration brought conformal coverage back to about 88--92\% on a new repository, but naive
thresholds still overshot. Prediction sets, which keep a human in the loop for ambiguous issues,
transfer better than a single auto-label threshold.

\paragraph{Deferral protects labels, not reporters.} Flips under rewording concentrate where the
service is unsure, so almost none are auto-applied. But a reporter who calls their post ``a silly
question'' is auto-labeled half as often and waits for a human instead. Measure workload and wait
time by reporter group, not only flips, and keep reporter self-descriptions out of the classifier's
input where possible.

\paragraph{Show evidence as context, not as cause.} The LLM's quotes are genuine, but we could not
show that they matter more than other text; a triage UI should present them as relevant passages,
not as the reason for the label.

\paragraph{Monitor inputs, not only the model.} A nearest-neighbor novelty score over the labeled
history caught a new repository within 13 requests at under 1\% false alarms; the model's confidence
did not, and disagreement between models alarmed falsely in 5--7\% of clean windows.

\section{Threats to Validity}\label{sec:threats}

\paragraph{Construct.} Labels are the repositories' own type labels, which mix issue content with
project policy (React routes questions elsewhere). The 2025--26 set maps each repository's label names
to the three types; issues with several type labels are excluded. Rewrites are synthetic. The
non-native rewrite was validated automatically (code blocks restored, length and technical tokens
preserved; 161 of 210 valid) and occasionally alters template fields; a human validation of
meaning preservation on a sample (\texttt{report/audit/rewrite\_audit.csv}) is pending. Removing code
is treated as reduced input, not as a meaning-preserving change. The cue sentences are one wording
each; other phrasings of the same status may behave differently.

\paragraph{Internal.} The calibration issues come from the benchmark's training half and the test set
from its test half, as in the competition. Temperature, threshold candidates and naive thresholds are
fitted on one third of the calibration issues and certified on the other two thirds; nothing is tuned
on test data. All classifiers see the same 1{,}500 characters of each issue. Probabilities are read
from the first output token; a different prompt or label wording would change them.

\paragraph{External.} One LLM, five large repositories, three labels. Smaller projects, other label
sets, and stronger models (including fine-tuned encoders such as the SetFit competition baseline)
may behave differently. The monitoring results come from simulated streams mixed from our pools,
not from a live deployment. All runs used one laptop.

\paragraph{Conclusion validity.} Rewording tests use 77--210 variants per kind; differences of a few
points between classifiers are within noise, while the ``silly question'' effect on the few-shot LLM
is not (paired bootstrap CI [10.0, 20.0]). New-repository results for risk-controlled thresholds rest
on very few auto-labeled issues (TF-IDF: 3 of 750 at the 10\% target) and support no conclusion
beyond the low coverage. The explanation test is underpowered ($p{=}0.115$).

\section{Conclusion}
Treating issue triage as a deployed, selective service changes what matters. The LLM was less
accurate than a simple classifier on the benchmark but held up a year later and on a new repository.
With an honest protocol, a few hundred calibration issues certify a 10\% error rate for almost no
issues, and no certificate carried over to a new repository. A reporter's description of their own
post, not their newcomer status, moved the LLM toward \emph{question}; deferral kept those flips from
being applied but shifted the work to maintainers. The cheapest monitor, input novelty, was the most
useful. The harness, data collection scripts, recorded model scores and regression suite are
available for replication.
